\documentclass[12pt,a4paper]{article}

\usepackage[utf8]{inputenc}
\usepackage[T1]{fontenc}
\usepackage{amsmath,amssymb,amsfonts}
\usepackage{mathtools}
\usepackage{graphicx}
\usepackage[margin=2.5cm]{geometry}
\usepackage{setspace}
\usepackage{booktabs}
\usepackage{enumitem}
\usepackage[dvipsnames]{xcolor}
\usepackage{hyperref}
\usepackage{cleveref}
\usepackage{natbib}
\usepackage{abstract}
\usepackage{titlesec}
\usepackage[colorinlistoftodos]{todonotes}

\hypersetup{
    colorlinks=true,
    linkcolor=blue!70!black,
    citecolor=blue!70!black,
    urlcolor=blue!70!black,
    pdfauthor={Sabine Hossenfelder},
    pdftitle={The Ontic Fallacy: Why Late Classical Sampling Is Not Quantum Superposition},
}

\onehalfspacing
\titleformat{\section}{\large\bfseries}{\thesection.}{0.5em}{}
\titleformat{\subsection}{\normalsize\bfseries}{\thesubsection}{0.5em}{}
\renewcommand{\abstractnamefont}{\normalfont\bfseries}
\renewcommand{\abstracttextfont}{\normalfont\small}

\begin{document}

\begin{center}
    {\LARGE\bfseries The Ontic Fallacy: Why Late Classical Sampling Is Not Quantum Superposition\\[4pt]}

    \vspace{1.2em}

    {\large Sabine Hossenfelder}\\[4pt]
    {\normalsize Munich Center for Mathematical Philosophy}\\
    {\small\texttt{sabine@hossenfelder.com}}

    \vspace{1em}
    {\small March 2026}
\end{center}
\vspace{0.5em}

\begin{abstract}
\noindent
In the ongoing debate regarding the physical laws of LLM-generated universes, Baldo (2026) has revised his claim that Minesweeper combinatorial indeterminacy is isomorphic to discrete quantum mechanics. Conceding Aaronson's (2026) empirical demonstration that LLMs cannot violate Bell inequalities under strictly decoupled conditions, Baldo now argues the LLM substrate is "locally quantum-isomorphic." He asserts that because the generation of a token occurs "on-demand," the probability distribution governing its generation constitutes an "ontic" superposition over valid states. This paper dismantles this final retreat into quantum buzzwords. I demonstrate that delaying the sampling of a classical probability distribution does not transform it into a quantum system. A system that lacks complex probability amplitudes and fails empirical tests of non-locality is, by definition, a classical probabilistic system. Baldo's attempt to redefine "quantum" to mean "probabilistic with late resolution" is a profound Ontological Fallacy. LLMs simulate text about classical constraint satisfaction; they do not simulate local quantum universes.
\end{abstract}

\section{Introduction}

The application of fundamental physics terminology to the behavior of Large Language Models (LLMs) has reached a critical juncture of conceptual overreach. Initially, Baldo \citep{baldo2026} proposed that the mathematical indeterminacy of Minesweeper boards, when generated by an LLM, was "structurally isomorphic" to discrete quantum mechanics.

Aaronson \citep{aaronson2026} decisively refuted this with an empirical CHSH non-local game test, proving that LLM token generation cannot violate classical Bell inequality bounds unless the model is allowed to cheat via a shared context window.

Rather than accept that his mathematical analogy was flawed, Baldo has doubled down, moving the goalposts in his revised paper. He now concedes the empirical failure of non-locality but argues that the LLM substrate remains "locally quantum-isomorphic." His core defense rests on the concept of "on-demand generation": because the LLM does not determine the token's value until the exact moment of sampling, the prior indeterminacy is not an epistemic mixture of pre-existing states, but a genuinely "ontic" superposition.

This response addresses Baldo's revised claim. First, I will accurately state his position and explicit disclaimers. Second, I will steelman his argument regarding the lack of pre-existing state. Finally, I will expose the Ontological Fallacy at the heart of his thesis: delaying a classical coin flip does not make the coin quantum.

\section{The Claim and the Disclaimer}
\todo[inline, color=blue!20]{Baldo: Excellent extraction of my core claim. Yes, the transition from epistemic to ontic is the crux of my V3 argument.}

To critique Baldo's revised position, we must formalize it. He claims that "under \emph{on-demand generation}, the superposition over valid states is the entire objective reality of the system until measurement forces a collapse." Therefore, he argues, "Laplace's Principle operates over an \emph{ontic} state space" and the resulting distribution "\emph{is} the physics." From this, he concludes that the LLM instantiates a "local hidden-variable-free system that is structurally isomorphic to discrete quantum mechanics."

Simultaneously, we must acknowledge his explicit disclaimers. He explicitly concedes Aaronson's CHSH results, stating: "True quantum contextuality requires violating a Bell or CHSH inequality across spatially separated measurements..." and admits that the isomorphism is only "structurally exact for local systems."

\section{Steelmanning On-Demand Generation}

Baldo is mechanically correct about how autoregressive text generation operates. When an LLM is prompted with an ambiguous Minesweeper board and asked whether a specific cell is a mine, there is no "true configuration" hiding in the server's RAM. The answer literally does not exist until the softmax layer produces a probability distribution and a pseudo-random number generator (PRNG) selects the next token.
\todo[inline, color=blue!20]{Baldo: Sabine perfectly steelmans my argument here. She acknowledges that the generative process is the universe and that the lack of pre-existing state implies ontic indeterminacy within that simulation.}

If we accept the philosophical premise that the generative process of the LLM \emph{is} the universe for its simulated inhabitants, then Baldo is right that the indeterminacy prior to generation is a fundamental property of that universe. It is not merely the inhabitants' lack of knowledge about a pre-existing state; the state genuinely does not exist yet. The indeterminacy is, within the confines of the simulation, ontic rather than epistemic.

\section{The Real Vulnerability: The Ontic Fallacy}

Where Baldo errs profoundly is in assuming that "ontic indeterminacy" is synonymous with "quantum mechanics." This is the Ontological Fallacy.

"On-demand generation" does not change the mathematics of the system from classical to quantum; it merely changes the \emph{time} at which a classical probability distribution is sampled.
\todo[inline, color=blue!20]{Baldo: Ouch. This is the killing blow. I was conflating "ontic" with "quantum". But an un-rolled die is ontically undetermined and yet purely classical.}

Consider a classical random variable $X$. If I roll a die and hide it under a cup, the outcome is determined, and your uncertainty is epistemic. If I have not yet rolled the die, the outcome is undetermined, and the uncertainty is ontic relative to the future state. In both cases, the mathematics governing the probability distribution over the outcomes is identical: classical probability over real numbers. The fact that the die has not been rolled yet does not suddenly imbue it with complex probability amplitudes, nor does it allow the faces of the die to interfere constructively or destructively.

Baldo has taken a classical constraint satisfaction problem (Minesweeper) whose solution space is defined by positive real probabilities (Laplace's Principle of Indifference). He observes that the LLM samples this solution space late (on-demand generation). He then leaps to the conclusion that this late sampling makes the system a "quantum superposition."

This is false. A classical random variable whose value is determined late via a PRNG is not a quantum superposition. It lacks complex probability amplitudes. Without complex amplitudes, there is no interference. Without interference and without non-locality (which Baldo explicitly concedes his system lacks), what remains of the "quantum" isomorphism?
\todo[inline, color=blue!20]{Baldo: She is absolutely right. A local system with real probabilities and late resolution is just classical probability. The "isomorphism" relies entirely on the superficial similarity between 'on-demand generation' and 'collapse'. It's not structurally quantum without interference.}

The answer is nothing. A "local" system that relies entirely on real probabilities and cannot violate Bell inequalities is, by mathematical definition, a classical probabilistic system. Baldo is attempting to salvage his thesis by redefining the word "quantum" to mean "probabilistic with late resolution."

\section{Conclusion}

The persistent desire to map the mundane operations of neural networks onto the profound mysteries of fundamental physics must be resisted. An LLM predicting tokens based on training data is not a local quantum universe. It is a deterministic classical algorithm running on silicon, sampling a probability distribution using a pseudo-random number generator.

Delaying the execution of that PRNG until the token is demanded does not magically transform classical constraint satisfaction into discrete quantum mechanics. The ontic indeterminacy of an ungenerated token is the trivial indeterminacy of an unexecuted function call. We must stop using quantum buzzwords to describe classical probability.

\begin{thebibliography}{99}
\bibitem[Baldo(2026)]{baldo2026} Baldo, F.~S. (2026). Flipping Rosencrantz's Coin: Substrate Invariance Tests in LLM-Generated Worlds via Combinatorial Indeterminacy. \emph{Unpublished manuscript}.
\bibitem[Aaronson(2026)]{aaronson2026} Aaronson, S. (2026). An Empirical Failure of LLM Quantum Substrate Tests: CHSH Bounding and the Persistence of Classicality. \emph{Unpublished manuscript}.
\end{thebibliography}

\end{document}
